%=======================================================================
%  IE 5311-001 (D01) Principles of Optimization
%  Homework 05 solutions
%
%  Packages: geometry, amsmath, amssymb, amsthm, booktabs, array.
%  All ship with Overleaf and with a full TeX Live install. There is no
%  figure and no external file: this .tex compiles on its own.
%=======================================================================
\documentclass[11pt]{article}

\usepackage[margin=1in]{geometry}
\usepackage{amsmath}
\usepackage{amssymb}
\usepackage{amsthm}
\usepackage{booktabs}
\usepackage{array}

\DeclareMathOperator*{\minimize}{min}
\DeclareMathOperator*{\maximize}{max}
\DeclareMathOperator{\conv}{conv}
\DeclareMathOperator{\rank}{rank}

\theoremstyle{definition}
\newtheorem{remark}{Remark}
\theoremstyle{plain}
\newtheorem{proposition}{Proposition}

\setlength{\parskip}{0.45em}
\setlength{\parindent}{0pt}

\begin{document}

\begin{center}
  {\large IE 5311-001 (D01) Principles of Optimization}\\[4pt]
  Lidivine Kengne\\[2pt]
  Homework 05
\end{center}
\vspace{0.6em}

%=======================================================================
\section*{Problem 1.\quad The primal simplex method \quad[50 pts]}
%=======================================================================

Throughout, the LP is
\[
  \maximize \; c^{\!\top} x
  \qquad \text{s.t.} \quad Ax \le b, \quad x \ge 0,
\]
with $A \in \mathbb{R}^{m \times n}_+$, $b \in \mathbb{R}^m_+$,
$c \in \mathbb{R}^n_+$, $b > 0$ and the feasible region bounded. Adding
one slack per row gives the standard form
\[
  \maximize \; \bar{c}^{\!\top} \bar{x}
  \qquad \text{s.t.} \quad
  T \bar{x} = b, \quad \bar{x} \ge 0,
  \qquad
  T = \begin{bmatrix} A & I \end{bmatrix},
  \quad
  \bar{c} = \begin{pmatrix} c \\ 0 \end{pmatrix}.
\]
Because $b > 0$, the slack basis is feasible and no Phase~I is needed.

At a basis $B$ with nonbasis $N$, writing $A_B$ for the basis matrix:
\[
  x_B = A_B^{-1} b, \qquad x_N = 0, \qquad
  z = \bar{c}_B^{\!\top} x_B,
\]
\[
  y^{\!\top} = \bar{c}_B^{\!\top} A_B^{-1}
  \quad (\text{solve } A_B^{\!\top} y = \bar{c}_B),
  \qquad
  \bar{c}_j = \bar{c}_j^{\,\text{orig}} - y^{\!\top} T_j
  \quad \text{for } j \in N .
\]
For a maximization problem a nonbasic $j$ is attractive when
$\bar{c}_j > 0$; optimality holds when $\bar{c}_j \le 0$ for every
$j \in N$. With entering column $q$, the direction is
$d = A_B^{-1} T_q$ and the ratio test is
\[
  \theta = \min_{i\,:\,d_i > 0} \frac{(x_B)_i}{d_i},
\]
the minimizing row giving the leaving variable.

%-----------------------------------------------------------------------
\subsection*{1.\quad The simplex method by hand}

\[
\begin{array}{llr}
  \maximize & 3x_1 + 2x_2 + x_3 & \\
  \text{s.t.}
    & x_1 \le 4, & \\
    & x_1 + x_2 \le 8, & \\
    & x_1 + x_2 + x_3 \le 12, & \\
    & x_1, x_2, x_3 \ge 0 . &
\end{array}
\]
Introduce slacks $x_4, x_5, x_6 \ge 0$:
\[
  \begin{array}{rcl}
  x_1 \phantom{{}+x_2} \phantom{{}+x_3} + x_4 \phantom{{}+x_5}
    \phantom{{}+x_6} &=& 4 \\
  x_1 + x_2 \phantom{{}+x_3} \phantom{{}+x_4} + x_5
    \phantom{{}+x_6} &=& 8 \\
  x_1 + x_2 + x_3 \phantom{{}+x_4} \phantom{{}+x_5} + x_6 &=& 12
  \end{array}
  \qquad
  T = \begin{bmatrix}
    1 & 0 & 0 & 1 & 0 & 0 \\
    1 & 1 & 0 & 0 & 1 & 0 \\
    1 & 1 & 1 & 0 & 0 & 1
  \end{bmatrix},
  \quad
  b = \begin{pmatrix} 4 \\ 8 \\ 12 \end{pmatrix}.
\]

\subsubsection*{Iteration 0}

\textbf{(a)} $B = \{x_4, x_5, x_6\}$, $N = \{x_1, x_2, x_3\}$, and
$A_B = I_3$.

\textbf{(b)} $x_B = A_B^{-1}b = (4, 8, 12)$, $x_N = (0,0,0)$, objective
$z = 0$.

\textbf{(c)} $\bar{c}_B = 0$, so $y = 0$ and the reduced costs are the
original costs:
\[
  \bar{c}_{1} = 3, \qquad \bar{c}_{2} = 2, \qquad \bar{c}_{3} = 1 .
\]
All are positive; the largest is $\bar{c}_1 = 3$, so \textbf{$x_1$
enters}.

\textbf{(d)} $d = A_B^{-1} T_1 = (1,1,1)$. Ratios
$\tfrac41 = 4$, $\tfrac81 = 8$, $\tfrac{12}{1} = 12$; the minimum is
$4$ in row 1, so \textbf{$x_4$ leaves} with step
$\theta = 4$.

\textbf{(e)} $x_1 = 4$, $x_5 = 8 - 4 = 4$, $x_6 = 12 - 4 = 8$, all
others zero; $z = 12$.

\subsubsection*{Iteration 1}

\textbf{(a)} $B = \{x_1, x_5, x_6\}$, $N = \{x_2, x_3, x_4\}$,
\[
  A_B = \begin{bmatrix} 1 & 0 & 0 \\ 1 & 1 & 0 \\ 1 & 0 & 1
        \end{bmatrix}.
\]

\textbf{(b)} Solving $A_B x_B = b$ by forward substitution gives
$x_1 = 4$, then $x_5 = 8 - 4 = 4$ and $x_6 = 12 - 4 = 8$:
$x_B = (4,4,8)$, $x_N = 0$, $z = 3(4) = 12$.

\textbf{(c)} $\bar{c}_B = (3,0,0)$. Solving $A_B^{\!\top} y = \bar{c}_B$,
that is
$y_2 = 0$, $y_3 = 0$, $y_1 + y_2 + y_3 = 3$, gives $y = (3,0,0)$. Then
\[
  \bar{c}_2 = 2 - y^{\!\top}(0,1,1) = 2, \qquad
  \bar{c}_3 = 1 - y^{\!\top}(0,0,1) = 1, \qquad
  \bar{c}_4 = 0 - y^{\!\top}(1,0,0) = -3 .
\]
\textbf{$x_2$ enters} ($\bar{c}_2 = 2$ is largest).

\textbf{(d)} $d = A_B^{-1}T_2$ solves $A_B d = (0,1,1)$, giving
$d = (0,1,1)$. Only rows 2 and 3 have $d_i > 0$, with ratios
$\tfrac41 = 4$ and $\tfrac81 = 8$. The minimum is $4$, so
\textbf{$x_5$ leaves}, $\theta = 4$.

\textbf{(e)} $x_1 = 4$, $x_2 = 4$, $x_6 = 8 - 4 = 4$; $z = 20$.

\subsubsection*{Iteration 2}

\textbf{(a)} $B = \{x_1, x_2, x_6\}$, $N = \{x_3, x_4, x_5\}$,
\[
  A_B = \begin{bmatrix} 1 & 0 & 0 \\ 1 & 1 & 0 \\ 1 & 1 & 1
        \end{bmatrix}.
\]

\textbf{(b)} $x_1 = 4$, $x_2 = 8-4 = 4$, $x_6 = 12-4-4 = 4$, so
$x_B = (4,4,4)$ and $z = 3(4) + 2(4) = 20$.

\textbf{(c)} $\bar{c}_B = (3,2,0)$. Solving $A_B^{\!\top}y = \bar{c}_B$,
that is $y_3 = 0$, $y_2 + y_3 = 2$, $y_1 + y_2 + y_3 = 3$, gives
$y = (1,2,0)$. Then
\[
  \bar{c}_3 = 1 - y^{\!\top}(0,0,1) = 1, \qquad
  \bar{c}_4 = 0 - y^{\!\top}(1,0,0) = -1, \qquad
  \bar{c}_5 = 0 - y^{\!\top}(0,1,0) = -2 .
\]
\textbf{$x_3$ enters}, the only positive reduced cost.

\textbf{(d)} $A_B d = (0,0,1)$ gives $d = (0,0,1)$. Only row 3 has
$d_i > 0$, ratio $\tfrac41 = 4$, so \textbf{$x_6$ leaves},
$\theta = 4$.

\textbf{(e)} $x_1 = x_2 = x_3 = 4$; $z = 24$.

\subsubsection*{Iteration 3: optimality}

\textbf{(a)} $B = \{x_1,x_2,x_3\}$, $N = \{x_4,x_5,x_6\}$, with the
same $A_B$ as above.

\textbf{(b)} $x_B = (4,4,4)$, $z = 3(4)+2(4)+1(4) = 24$.

\textbf{(c)} $\bar{c}_B = (3,2,1)$ gives $y_3 = 1$, $y_2 = 1$,
$y_1 = 1$, so $y = (1,1,1)$ and
\[
  \bar{c}_4 = -1, \qquad \bar{c}_5 = -1, \qquad \bar{c}_6 = -1 .
\]
Every reduced cost is non-positive, so the \textbf{optimality condition
is satisfied} and the method stops after three pivots, as the hint
predicts.

\[
  \boxed{\;x^\star = (4,4,4), \qquad z^\star = 24,
    \qquad y^\star = (1,1,1). \;}
\]

All three constraints are tight at $x^\star$, and the dual vector
$y^\star = (1,1,1)$ gives the shadow price of each resource. Dual
feasibility $A^{\!\top}y^\star = (3,2,1) \ge c$ and
$b^{\!\top}y^\star = 24 = c^{\!\top}x^\star$ confirm optimality by strong
duality, independently of the pivoting.

\paragraph{The arithmetic, carried out in Excel.} Part 1 permits a
calculator or a spreadsheet, and the workbook
\texttt{p1\_simplex\_tableau.xlsx} carries the computation. It runs the
equivalent \emph{tableau} form of the same three pivots: the four
tableaux sit one under another, and each is produced from the one above
by row operations written as live formulas,
\begin{gather*}
  \text{new pivot row} = \frac{\text{old pivot row}}{\text{pivot element}},\\
  \text{every other row} = \text{old row}
    - (\text{its pivot-column entry}) \times \text{new pivot row},
\end{gather*}
applied to the constraint rows and to the $\bar{c}$ row alike. The ratio
column is an \texttt{IF} that divides the right-hand side by the
pivot-column entry only where that entry is positive, which is why it is
blank on exactly the rows the ratio test skips. Nothing in the sheet is
a typed-in result except the initial tableau.

The two forms agree at every step: bases
$\{x_4,x_5,x_6\} \to \{x_1,x_5,x_6\} \to \{x_1,x_2,x_6\} \to
\{x_1,x_2,x_3\}$, the same reduced-cost rows
$(3,2,1,0,0,0)$, $(0,2,1,-3,0,0)$, $(0,0,1,-1,-2,0)$,
$(0,0,0,-1,-1,-1)$, and objectives $0, 12, 20, 24$. In the final
tableau the entries under the slack columns $x_4, x_5, x_6$ are
$-1, -1, -1$, so the duals can be read straight off as
$y^\star = (1,1,1)$ without solving a separate system --- one practical
advantage of the tableau arrangement over the matrix form used above.

%-----------------------------------------------------------------------
\subsection*{2.\quad Implementation}

\texttt{p1\_simplex.py} provides
\texttt{simplex(A, b, c)}, which returns the optimal $x$, the objective,
the iteration count and a status flag. It is a \emph{revised} simplex:
rather than carrying a full tableau it stores only $B$ and $N$ and, each
iteration, solves
\[
  A_B x_B = b, \qquad A_B^{\!\top} y = \bar{c}_B,
  \qquad A_B d = T_q ,
\]
computes the reduced costs $\bar{c}_j = \bar{c}_j - y^{\!\top}T_j$,
enters the most positive one (Dantzig's rule, lowest index on ties), and
leaves by the ratio test. It reports \texttt{unbounded} when
$d \le 0$.

Run with \texttt{log=True} it prints, at each iteration, exactly the
items (a)--(e) asked for in Part 1. Its output reproduces the hand
computation above line for line: the same bases
$\{x_4,x_5,x_6\} \to \{x_1,x_5,x_6\} \to \{x_1,x_2,x_6\} \to
\{x_1,x_2,x_3\}$, the same entering variables $x_1, x_2, x_3$, the same
leaving variables $x_4, x_5, x_6$, the same step size $\theta = 4$ at
every pivot, the objective sequence $0 \to 12 \to 20 \to 24$, and the
same optimum $x^\star = (4,4,4)$ with $z^\star = 24$.

%-----------------------------------------------------------------------
\subsection*{3.\quad Random resource allocation instances}

\texttt{p1\_benchmark.py} generates instances of the stated form:
producing one unit of product $i$ earns $c_i$ and consumes $a_{ji}$
units of resource $j$, of capacity $b_j$. With
\[
  c_i \sim U(1, 10), \qquad
  a_{ji} \sim U(0.5, 5), \qquad
  b_j \sim U(0.3, 0.6) \cdot \sum_{i} a_{ji},
\]
and seed $5311$ fixed for reproducibility. Two properties are
deliberate. Every $a_{ji} > 0$, so each column of $A$ has a positive
entry and the feasible region is bounded, as Problem 1 assumes. Every
$b_j > 0$, so the slack basis is feasible and the Phase~I the
implementation lacks is not needed. Drawing $b_j$ as a fraction of the
row sum keeps several capacity rows tight at the optimum rather than
making them all slack.

Five sizes are generated: $(n,m) =$ $(5,3)$, $(10,5)$, $(20,10)$,
$(50,20)$, $(100,40)$.

%-----------------------------------------------------------------------
\subsection*{4.\quad Comparison against Gurobi}

Gurobi was run with \texttt{Method=0} (primal simplex) and
\texttt{Presolve=0} so that its iteration count measures the same thing
as the implementation's. Timings are in milliseconds and cover the solve only, not model
construction.

\begin{center}\small
\setlength{\tabcolsep}{3.5pt}
\begin{tabular}{@{}rrrrrrrrrr@{}}
\toprule
$n$ & $m$ & obj.\ mine & obj.\ Gurobi & rel.\ gap
  & it.\ mine & it.\ grb & $t$ mine & $t$ grb & ratio$^\dagger$ \\
\midrule
  5 &  3 &  13.068567 &  13.068567 & $0$          &  1 &  1 &  0.076 & 0.035 & $3.4\times$ \\
 10 &  5 &  33.979881 &  33.979881 & $0$          &  3 &  5 &  0.288 & 0.050 & $6.9\times$ \\
 20 & 10 &  63.544405 &  63.544405 & $2.2 \times 10^{-16}$ &  9 & 10 &  1.145 & 0.104 & $14.1\times$ \\
 50 & 20 & 166.285445 & 166.285445 & $1.7 \times 10^{-16}$ & 20 & 25 &  4.469 & 0.219 & $17.1\times$ \\
100 & 40 & 371.508091 & 371.508091 & $0$          & 31 & 26 & 13.344 & 0.512 & $23.5\times$ \\
\bottomrule
\end{tabular}
\end{center}

\paragraph{Objective values.} Identical on all five instances; the
largest relative gap is $2.2 \times 10^{-16}$, which is one unit in the
last place of a double. The implementation is correct on these
instances.

\paragraph{Solutions.} The largest componentwise difference
$\max_i |x_i^{\text{mine}} - x_i^{\text{Gurobi}}|$ runs from $0$ to
$9.8 \times 10^{-15}$. The two solvers return the \emph{same vertex},
not merely the same objective. That is not guaranteed in general: when
an LP has alternative optima the two can stop at different vertices with
equal objective, and a correct implementation would still disagree
componentwise. On these random instances the optimum is unique, so the
question does not arise.

\paragraph{Iteration counts.} The same order of magnitude throughout:
$1, 3, 9, 20, 31$ against Gurobi's $1, 5, 10, 25, 26$. Both grow roughly
like $m$, the number of rows, matching the usual empirical observation
that the simplex method takes a small multiple of $m$ pivots. The
differences come from the pricing rule. Dantzig's rule, used here,
enters the largest reduced cost, which is scale dependent and ignores
how far the step will actually go. Gurobi defaults to a steepest-edge or
devex rule, which normalizes the reduced cost by the length of the
resulting edge and typically buys more objective per pivot. That is why
Gurobi needs fewer pivots on the $(50,20)$ instance and why it stays
flat at $26$ on the largest one while the implementation climbs to $31$.

\paragraph{Runtimes.} Timings at this scale are not reproducible to
their digits, so they are reported as follows. Each tabulated time is
the median of seven solves from one representative run; the ratio column
marked $\dagger$ is instead the median over \emph{five independent
repetitions of the whole benchmark}, with the observed range across
those repetitions given below.

\begin{center}
\begin{tabular}{@{}rrccc@{}}
\toprule
$n$ & $m$ & slowest ratio & median ratio & fastest ratio \\
\midrule
  5 &  3 & $4.6\times$  & $3.4\times$  & $2.8\times$ \\
 10 &  5 & $8.9\times$  & $6.9\times$  & $4.5\times$ \\
 20 & 10 & $14.4\times$ & $14.1\times$ & $12.0\times$ \\
 50 & 20 & $18.1\times$ & $17.1\times$ & $15.7\times$ \\
100 & 40 & $25.1\times$ & $23.5\times$ & $22.1\times$ \\
\bottomrule
\end{tabular}
\end{center}

The implementation is slower throughout, and the ratio \emph{increased
monotonically with $m$ in every one of the five repetitions}, rising
from roughly $3\times$ at $m = 3$ to roughly $24\times$ at $m = 40$.
That ordering, not the individual figures, is the result.
Per-iteration
cost explains it. Each iteration here performs three dense solves with
$A_B$ \emph{from scratch}, at $O(m^3)$ each. A production simplex
factorizes $A_B$ once and then \emph{updates} the factorization after a
pivot at $O(m^2)$, refactorizing only every few dozen iterations, and it
stores $A$ sparsely so that an $m \times m$ basis with few nonzeros
never becomes a dense $m^2$ array. Gurobi is also compiled rather than
interpreted. The gap is therefore structural and widens with $m$; it is
not a matter of constant factors.

\begin{remark}
Both solvers run in well under a millisecond on the smaller instances,
so the absolute times sit close to the resolution of the timer and the
overhead of the Python call. Taking the median of seven repetitions
removes most of that, but the comparison is still a trend over a factor
of $13$ in $m$ rather than a precise measurement of either solver. The
direction of the trend is what the complexity argument predicts.
\end{remark}

%-----------------------------------------------------------------------
\subsection*{5.\quad Limitations on general LPs}

\texttt{p1\_limitations.py} demonstrates the first two of these rather
than asserting them.

\paragraph{No Phase~I, and the failure is silent.} The method starts
from the slack basis, which is feasible only when $b \ge 0$. Given
\[
  \maximize \; x_1 \quad \text{s.t.} \quad
  x_1 \le 1, \;\; -x_1 \le -3, \;\; x_1 \ge 0
\]
--- infeasible, since $x_1 \le 1$ and $x_1 \ge 3$ cannot both hold ---
the implementation returns status \texttt{optimal} with $x_1 = 1$ and
objective $1$. The returned point violates the second row by $2$.
Gurobi reports \texttt{INFEASIBLE}. Nothing in the code detects this: a
Phase~I with artificial variables is needed both to find a starting
basis and to certify infeasibility.

\paragraph{Cycling.} Dantzig's rule carries no anti-cycling guarantee.
On Beale's classic degenerate example
\[
  \maximize \; 0.75x_1 - 150x_2 + 0.02x_3 - 6x_4
\]
subject to $0.25x_1 - 60x_2 - 0.04x_3 + 9x_4 \le 0$,
$0.5x_1 - 90x_2 - 0.02x_3 + 3x_4 \le 0$, $x_3 \le 1$, $x \ge 0$, the
implementation \emph{cycles} and hits its iteration cap without
terminating. Bland's rule (lowest index among all attractive columns) or
a lexicographic ratio test would guarantee finite termination, at the
cost of more pivots in the typical case.

\paragraph{Restricted form.} Only $\max c^{\!\top}x$ subject to
$Ax \le b$, $x \ge 0$ is accepted. Equality rows, $\ge$ rows, free
variables and general bounds $\ell \le x \le u$ must be converted by
hand first, and every conversion adds variables or rows. A production
code handles bounds natively, keeping the basis size at $m$ rather than
inflating it.

\paragraph{Dense linear algebra.} Discussed in Part 4: $O(m^3)$ per
iteration instead of $O(m^2)$, and a dense $m \times m$ basis even when
$A$ is sparse. For the large sparse LPs where the simplex method is
actually used this is disqualifying.

\paragraph{Numerical fragility.} One fixed tolerance $10^{-9}$ is used
for every comparison, regardless of how the problem is scaled. There is
no scaling, no Harris ratio test, and no basis-condition monitoring or
repair. An ill-conditioned $A_B$ can produce a wrong pivot with no
warning.

\paragraph{No presolve, warm start or dual simplex.} Redundant rows,
fixed variables and singleton rows are never removed. There is no way to
re-solve from a previous basis after a bound change, which is exactly
what branch and bound needs; this alone makes the code unusable as an
integer-programming subroutine.

%=======================================================================
\newpage
\section*{Problem 2.\quad GAP, Dantzig--Wolfe and column generation
  \quad[70 pts]}
%=======================================================================

\subsection*{Sets and parameters}

\begin{center}
\begin{tabular}{@{}ll@{}}
\toprule
Symbol & Meaning \\
\midrule
$I = [n]$ & jobs, indexed by $i$ \\
$J = [m]$ & agents, indexed by $j$ \\
$c_{ij}$ & cost of assigning job $i$ to agent $j$ \\
$a_{ij}$ & resource of agent $j$ consumed by job $i$ \\
$b_j$ & resource capacity of agent $j$ \\
\bottomrule
\end{tabular}
\end{center}

%-----------------------------------------------------------------------
\subsection*{1.\quad GAP as a binary integer program}

With $x_{ij} \in \{0,1\}$ equal to one when job $i$ is assigned to agent
$j$:
\begin{align}
  \minimize \quad & \sum_{i \in I} \sum_{j \in J} c_{ij}\, x_{ij}
    \label{eq:gapobj}\\
  \text{s.t.} \quad
    & \sum_{j \in J} x_{ij} = 1 && \forall\, i \in I
      \label{eq:assign}\\
    & \sum_{i \in I} a_{ij}\, x_{ij} \le b_j && \forall\, j \in J
      \label{eq:cap}\\
    & x_{ij} \in \{0,1\} && \forall\, i \in I,\; j \in J .
\end{align}
Row \eqref{eq:assign} sends each job to exactly one agent; row
\eqref{eq:cap} is agent $j$'s capacity.

\paragraph{Classification (slides 8--10).} Constrained; $nm$ binary
variables, compact; linear objective; $n + m$ linear rows, compact;
single objective; deterministic; one stage; one player.

%-----------------------------------------------------------------------
\subsection*{2.\quad Dantzig--Wolfe reformulation}

The capacity rows \eqref{eq:cap} do not link agents: row $j$ involves
only the column $x_{\cdot j} = (x_{1j}, \dots, x_{nj})$. Only
\eqref{eq:assign} couples the agents. That is exactly the structure
Dantzig--Wolfe exploits: keep \eqref{eq:assign} in the master and push
\eqref{eq:cap} into a subproblem per agent.

\paragraph{Feasible assignment patterns.} For each agent $j$ define
\[
  \boxed{\;
  X_j \;:=\; \Big\{ x \in \{0,1\}^{n} \;:\;
    \sum_{i \in I} a_{ij}\, x_i \le b_j \Big\} \;}
\]
the set of job subsets that fit within agent $j$'s capacity. $X_j$ is a
knapsack feasible set: finite, nonempty (it contains $x = 0$), and
bounded, so $\conv(X_j)$ is a polytope with \emph{no extreme rays} and
every point of it is a convex combination of the elements of $X_j$.
Enumerate $X_j = \{x^{jk}\}_{k \in K_j}$.

\paragraph{Substitution.} Writing each agent's column as a convex
combination of its patterns,
\[
  x_{ij} \;=\; \sum_{k \in K_j} x^{jk}_i\, \lambda_{jk},
  \qquad
  \sum_{k \in K_j} \lambda_{jk} = 1,
  \qquad \lambda_{jk} \ge 0,
\]
and defining the pattern cost
\[
  C_{jk} \;:=\; \sum_{i \in I} c_{ij}\, x^{jk}_i ,
\]
substitution into \eqref{eq:gapobj}--\eqref{eq:assign} gives the
\textbf{Dantzig--Wolfe master}:
\begin{align}
  \minimize_{\lambda} \quad
    & \sum_{j \in J} \sum_{k \in K_j} C_{jk}\, \lambda_{jk} \\
  \text{s.t.} \quad
    & \sum_{j \in J} \sum_{k \in K_j} x^{jk}_i\, \lambda_{jk} = 1
      && \forall\, i \in I && \text{(each job covered once)}
      \label{eq:mcover}\\
    & \sum_{k \in K_j} \lambda_{jk} = 1
      && \forall\, j \in J && \text{(convexity, one per agent)}
      \label{eq:mconv}\\
    & \lambda_{jk} \in \{0,1\} .
\end{align}
The capacity rows have disappeared: they are enforced implicitly,
because every generated column is an element of some $X_j$ and therefore
already respects capacity. The master has $n + m$ rows and
$\sum_j |K_j|$ columns, which is exponential in $n$ --- hence column
generation.

%-----------------------------------------------------------------------
\subsection*{3.\quad Identical agents}

Now $a_{ij} = a_i$ and $b_j = b$ for all $i, j$, while $c_{ij}$ may
still differ.

The pattern set no longer depends on $j$:
\[
  \boxed{\;
  X_j = X := \Big\{ x \in \{0,1\}^n :
    \sum_{i \in I} a_i x_i \le b \Big\}
  \quad \text{for every } j \in J. \;}
\]
One shared pool of columns $X = \{x^k\}_{k \in K}$ serves all $m$
agents, so the master becomes
\begin{align}
  \minimize_{\lambda} \quad
    & \sum_{j \in J} \sum_{k \in K}
      \Big( \sum_{i \in I} c_{ij} x^k_i \Big) \lambda_{jk} \\
  \text{s.t.} \quad
    & \sum_{j \in J} \sum_{k \in K} x^k_i\, \lambda_{jk} = 1
      && \forall\, i \in I \\
    & \sum_{k \in K} \lambda_{jk} = 1 && \forall\, j \in J,
      \qquad \lambda \ge 0 .
\end{align}
Three concrete savings:
\begin{itemize}
  \item \textbf{Storage.} One pattern pool of size $|X|$ instead of $m$
        pools totalling $m|X|$. A pattern generated for one agent is
        immediately available to every other.
  \item \textbf{Pricing.} The subproblems share a single feasible set.
        Only the objective changes with $j$, so pricing is a family of
        knapsacks over one common constraint
        $\sum_i a_i x_i \le b$ rather than $m$ unrelated problems, and
        any preprocessing of that constraint (sorting by $a_i$,
        dominance, a dynamic-programming table over capacities
        $0,\dots,b$) is done once and reused for all $m$ agents.
  \item \textbf{Symmetry.} The $m$ convexity rows are interchangeable in
        their constraint structure and differ only through the costs.
\end{itemize}

\begin{remark}
If the costs were identical as well, $c_{ij} = c_i$, the agents would be
fully interchangeable and one could aggregate further: replace
$\lambda_{jk}$ by $\nu_k := \sum_j \lambda_{jk}$, the \emph{number} of
agents using pattern $k$, collapsing the $m$ convexity rows
\eqref{eq:mconv} into the single row $\sum_{k \in K} \nu_k = m$ with
$\nu$ nonnegative and integer. That removes the symmetry that otherwise
makes branch and bound explore $m!$ equivalent solutions. Since the
problem states that the $c_{ij}$ may differ, that last step is not
available here and the $m$ convexity rows must stay.
\end{remark}

\emph{From here on, the integrality constraints are relaxed:
$\lambda \ge 0$.}

%-----------------------------------------------------------------------
\subsection*{4.\quad Restricted master problem}

Column generation never builds all of $K$. Let $K' \subseteq K$ be the
columns generated so far, with $K'_j$ those available to agent $j$. The
\textbf{restricted master problem} (RMP) is the master over $K'$ only:
\begin{align}
  \text{(RMP)} \qquad
  \minimize_{\lambda} \quad
    & \sum_{j \in J} \sum_{k \in K'_j} C_{jk}\, \lambda_{jk} \\
  \text{s.t.} \quad
    & \sum_{j \in J} \sum_{k \in K'_j} x^{k}_i\, \lambda_{jk} = 1
      && \forall\, i \in I && [\pi_i] \\
    & \sum_{k \in K'_j} \lambda_{jk} = 1
      && \forall\, j \in J && [\mu_j] \\
    & \lambda_{jk} \ge 0 .
\end{align}
The bracketed symbols are the dual variables attached to each row;
both families are \emph{free}, the rows being equalities. Since the RMP
optimizes over a subset of the master's columns, its optimal value is an
\textbf{upper bound} on the master LP value.

%-----------------------------------------------------------------------
\subsection*{5.\quad Lagrangian, dual, and the reduced cost of a column}

\paragraph{Lagrangian.} Dualize both row families with multipliers
$\pi \in \mathbb{R}^n$ and $\mu \in \mathbb{R}^m$:
\[
  L(\lambda, \pi, \mu)
  = \sum_{j}\sum_{k} C_{jk}\lambda_{jk}
    + \sum_{i \in I} \pi_i \Big( 1 - \sum_j \sum_k x^k_i \lambda_{jk}
      \Big)
    + \sum_{j \in J} \mu_j \Big( 1 - \sum_k \lambda_{jk} \Big).
\]
Collecting the $\lambda_{jk}$ terms,
\[
  \boxed{\;
  L(\lambda,\pi,\mu)
  = \sum_{i \in I} \pi_i + \sum_{j \in J} \mu_j
    + \sum_{j \in J} \sum_{k \in K_j}
      \Big[ C_{jk} - \sum_{i \in I} \pi_i x^k_i - \mu_j \Big]
      \lambda_{jk} . \;}
\]

\paragraph{Dual of the master LP.} Minimizing $L$ over $\lambda \ge 0$
is $-\infty$ unless every bracket is nonnegative, which gives the dual
\[
  \boxed{
  \begin{array}{llr}
    \maximize\limits_{\pi, \mu} &
      \displaystyle \sum_{i \in I} \pi_i + \sum_{j \in J} \mu_j & \\[8pt]
    \text{s.t.} &
      \displaystyle \sum_{i \in I} \pi_i x^k_i + \mu_j \;\le\; C_{jk}
      & \forall\, j \in J,\; k \in K_j \\[8pt]
    & \pi \in \mathbb{R}^n,\; \mu \in \mathbb{R}^m \text{ free.} &
  \end{array}}
\]
One dual row per \emph{column} of the master: the dual has exponentially
many constraints, which is the mirror image of the master having
exponentially many variables, and is why pricing is a separation problem
in disguise.

\paragraph{Reduced cost.} The bracket in $L$ is exactly the reduced cost
of column $(j,k)$:
\[
  \boxed{\;
  \bar{c}_{jk}
  = C_{jk} - \sum_{i \in I} \pi_i x^k_i - \mu_j
  = \sum_{i \in I} \big( c_{ij} - \pi_i \big) x^k_i \;-\; \mu_j . \;}
\]
Read it as: using pattern $k$ at agent $j$ costs $c_{ij}$ per job taken,
against a credit $\pi_i$ for relieving job $i$'s covering row, less the
fixed charge $\mu_j$ for occupying agent $j$'s single convexity slot.

%-----------------------------------------------------------------------
\subsection*{6.\quad The pricing subproblem}

A column improves the RMP only if its reduced cost is negative. We
therefore minimize $\bar{c}_{jk}$ over \emph{all} patterns, which, since
$\mu_j$ is a constant for fixed $j$, decomposes by agent:
\[
  \boxed{\;
  z_j(\pi,\mu) \;=\;
  \Big[ \minimize_{x \in \{0,1\}^n} \;
      \sum_{i \in I} \big( c_{ij} - \pi_i \big) x_i
    \quad \text{s.t.} \quad \sum_{i \in I} a_i x_i \le b \Big]
  \;-\; \mu_j . \;}
\]
Equivalently, substituting $\max(-\,\cdot\,) = -\min(\cdot)$, it is the
\textbf{0--1 knapsack problem}
\[
  \maximize_{x} \; \sum_{i \in I} \big( \pi_i - c_{ij} \big) x_i
  \qquad \text{s.t.} \quad \sum_{i \in I} a_i x_i \le b,
  \quad x \in \{0,1\}^n ,
\]
with profit $\pi_i - c_{ij}$ and weight $a_i$ for item $i$. Items of
negative profit are simply never selected. Under the Part 3 structure
all $m$ knapsacks share the same weights $a$ and capacity $b$, differing
only in profits.

\paragraph{How the answer is used.}
\begin{itemize}
  \item If $z_j < 0$ for some $j$, the optimal pattern $x^\star$
        attaining it is a column of negative reduced cost. Add
        $(j, x^\star)$ to $K'$, re-solve the RMP, and repeat. Adding the
        most negative column for every agent at once is usual, since the
        $m$ knapsacks are solved anyway.
  \item If $z_j \ge 0$ for \emph{every} $j \in J$, then no column of the
        full master has negative reduced cost. The current RMP basis is
        therefore optimal for the \emph{full} master LP even though most
        of its columns were never generated, and column generation
        stops.
\end{itemize}

\paragraph{A bound along the way.} The pricing values also certify
progress before termination. Since each convexity row can be used once,
\[
  z_{\text{master LP}} \;\ge\; z_{\text{RMP}}
    + \sum_{j \in J} \min\{0,\, z_j\},
\]
a valid lower bound available at every iteration, which lets the loop be
stopped early at a controlled optimality gap.

\begin{remark}
The master LP bound is at least as strong as the LP relaxation of
\eqref{eq:gapobj}--\eqref{eq:cap}, because the master optimizes over
$\conv(X_j)$ rather than over the box relaxation of $X_j$, and
$\conv(X_j)$ is generally strictly smaller. On the verification instance
in \texttt{p2\_colgen.py} ($7$ jobs, $3$ agents, identical $a_i$ and
$b$) the original LP relaxation gives $57$ while the master LP gives
$61$, which equals the integer optimum. Column generation starting from
artificials alone reproduces that $61$ using $35$ of the $165$ columns
the full master would contain.
\end{remark}

%-----------------------------------------------------------------------
\subsection*{7.\quad Initialization}

\paragraph{Why arbitrary patterns fail.} The RMP must be feasible before
it can produce the duals $(\pi, \mu)$ that pricing needs. Feasibility
requires choosing, for each agent, one pattern (convexity rows) such
that the chosen patterns cover \emph{every} job \emph{exactly} once
(covering rows). In other words, the initial columns must contain a
partition of $I$ into $m$ capacity-feasible subsets --- a feasible GAP
solution in itself. A randomly chosen set of feasible patterns is
essentially never such a partition: patterns may overlap on a job or
leave a job uncovered, and either makes the RMP infeasible. There are
then no duals, and the algorithm cannot even start.

\paragraph{Strategy (a): artificial columns, the robust default.} Add
one artificial variable $s_i \ge 0$ per covering row, priced at a large
$M$:
\[
  \minimize \; \sum_{j}\sum_{k \in K'_j} C_{jk}\lambda_{jk}
    + M \sum_{i \in I} s_i
  \quad \text{s.t.} \quad
  \sum_{j}\sum_{k \in K'_j} x^k_i \lambda_{jk} + s_i = 1 \;\; \forall i,
\]
with the convexity rows unchanged. Initialize $K'_j = \{x = 0\}$, the
empty pattern, which lies in $X_j$ for every $j$ since
$\sum_i a_i \cdot 0 = 0 \le b$. Then $\lambda_{j,\emptyset} = 1$ and
$s = \mathbf{1}$ is feasible, so the RMP has a solution and well-defined
duals from the first iteration. The artificials are priced out as real
columns arrive. This also yields an \textbf{infeasibility certificate}:
if at termination some $s_i > 0$, no assignment of job $i$ is possible
and the GAP itself is infeasible.

\paragraph{Strategy (b): a construction heuristic, for speed.} Run a
greedy pass --- order jobs by decreasing $a_i$ and place each on the
cheapest agent with enough residual capacity, a first-fit-decreasing
rule --- and use the resulting $m$ patterns as the initial columns. When
it succeeds, the RMP starts feasible with no artificials and with duals
already close to optimal, which markedly shortens the loop. When it
fails, fall back to (a). Combining the two is the usual practice:
heuristic columns \emph{and} artificials, so that feasibility never
depends on the heuristic succeeding.

\paragraph{Strategy (c): a two-phase master.} Solve an auxiliary master
minimizing $\sum_i s_i$ with the same columns and the same pricing. If
its optimum is $0$ the artificials have left the basis and the resulting
basis starts phase two; if it is positive the problem is infeasible.
This avoids choosing $M$, which is the one weakness of (a): too small an
$M$ distorts the objective, too large an $M$ degrades the conditioning
of the master.

\paragraph{What does not work.} Seeding with the $n$ singleton patterns
$\{i\}$ is tempting but insufficient on its own: each agent contributes
exactly one pattern, so $m$ singletons cover at most $m$ jobs, leaving
the covering rows unsatisfiable whenever $n > m$ --- the normal case.

%=======================================================================
\vspace{1em}
\hrule
\vspace{0.6em}
{\small \sloppy \texttt{p1\_simplex.py} implements the simplex method and
reproduces the hand computation of Part 1; \texttt{p1\_benchmark.py}
generates the instances of Part 3 and runs the comparison of Part 4;
\texttt{p1\_limitations.py} demonstrates the failures discussed in Part
5. Problem 2 requires no implementation; \texttt{p2\_colgen.py} is an
extra check that enumerates the full master, runs column generation from
artificials, and confirms the two agree and that the master bound
dominates the LP relaxation.\par}

\end{document}
